\uuid{6MeA}
\titre{Changement de variables}
\auteur{hanani ; mzari}
\organisation{exo7}
\licence{CC BY-NC-SA 4.0}
\datecreate{2020-01-20}
\niveau{L1}
\module{Analyse}
\chapitre{Calcul d'intégrales}
\sousChapitre{Primitives diverses}
\difficulte{3}
\theme{Primitives, changement de variables}
\exo7qcmid{lille-2.primitives.036}

\contenu{
\begin{question}
\qclassification{127.11}
\qtags{motcle=changement de variables}
Soit $\displaystyle f(x)=\frac{1}{\mathrm{e}^x+\mathrm{e}^{-x}}$. Parmi les affirmations suivantes, cocher celles qui sont vraies :
\begin{answers}
\good{$f$ admet une primitive sur $\Rr$.}
\bad{La fonction $F$ telle que $\displaystyle F(x)=\ln \left(\mathrm{e}^x+\mathrm{e}^{-x}\right)$ est une primitive de $f$ sur $\Rr$.}
\good{La fonction $F$ telle que $\displaystyle F(x)=\frac{\pi}{4}-\arctan \left(\mathrm{e}^{-x}\right)$ est une primitive de $f$ sur $\Rr$.}
\bad{La primitive de $f$ sur $\Rr$ qui s'annule en $0$ est d\'efinie par $\displaystyle F(x)=\frac{\pi}{4}-\arctan \left(\mathrm{e}^{x}\right)$.}
\end{answers}
\begin{explanations}
La fonction $f$ est continue sur $\Rr$, donc elle y admet des primitives. Avec $t=\mathrm{e}^{-x}$, on a : $\mathrm{d}t=-\mathrm{e}^{-x}\, \mathrm{d}x$ et
$$\int \frac{\mathrm{d}x}{\mathrm{e}^x+\mathrm{e}^{-x}}=\int \frac{\mathrm{e}^{-x}\, \mathrm{d}x}{1+\mathrm{e}^{-2x}}=-\int \frac{\mathrm{d}t}{1+t^2}=-\arctan (t)+k=-\arctan \left(\mathrm{e}^{-x}\right)+k,\; k\in \Rr.$$
La condition $F(0)=0$ implique que $\displaystyle k=\frac{\pi}{4}$.
\end{explanations}
\end{question}
}
